\chapter{External Distribution Sort, Over-Sampling, and Random Sampling Foundations}
\label{chap:external-distribution-sort}

\begin{center}
  \textbf{Lecture 7}\\
  \emph{External Distribution Sort, Over-Sampling Technique, and Foundations of Random Sampling}\\[0.5em]
  Date: September 30, 2026
\end{center}

\section{The External Sorting Problem and Multi-Way Partitioning}
\label{sec:external-sorting-multiway}

\subsection{The Computational I/O Model (Secondary Memory)}
In systems processing massive datasets, the size of the dataset to be sorted $S$ far exceeds the capacity of random access internal memory (RAM). We adopt the standard two-level secondary memory model introduced by Aggarwal and Vitter~\cite{aggarwal1988}:
\begin{itemize}
  \item $N$ (or $n$): total number of records in sequence $S = \langle S[1], \dots, S[n] \rangle$;
  \item $M$: internal memory (RAM) capacity, measured in number of records ($M \ll n$);
  \item $B$: disk transfer block size, measured in number of records transferred per single Input/Output (I/O) operation.
\end{itemize}

The number of blocks that can reside simultaneously in internal memory is $\frac{M}{B}$.

To sort the dataset $S$ with optimal I/O complexity, a standard binary recursive partition algorithm (such as classical Quicksort) is insufficient, as it yields a recursion tree of depth $\Theta(\log_2(n/B))$, which is asymptotically too deep and incurs excessive random disk seeks.

Instead, algorithms rely on multi-way partitioning (\textbf{Multi-Way Partitioning} or \textbf{External Distribution Sort}), dividing the dataset recursively into $k$ disjoint partitions (called \textit{buckets}), where the partition fan-out is calibrated to available memory:
\begin{equation}
  k = \mathcal{O}\left(\frac{M}{B}\right).
\end{equation}

\subsection{Formal Definition of Buckets and Pivots}
Given the unsorted array $S[1 \dots n]$, the goal is to select an ordered set of $k - 1$ sentinel elements, called \textbf{pivots}:
\begin{equation}
  s_1 \le s_2 \le \dots \le s_{k-1}.
\end{equation}
The notation $s$ emphasizes that these elements originate from a sampling procedure (\textit{sampled}).

For analytical convenience and to ensure boundary completeness, we introduce two conceptual sentinel pivots:
\begin{equation}
  s_0 \coloneqq -\infty, \qquad s_k \coloneqq +\infty.
\end{equation}

\begin{definition}[Buckets $B_i$]
The $k$ buckets $B_1, B_2, \dots, B_k$ are defined as the disjoint partitions formed by all elements of $S$ that fall within the respective half-open intervals determined by adjacent pivots:
\begin{equation}
  B_i \coloneqq \Big\{ S[j] \in S \;\mid\; s_{i-1} < S[j] \le s_i \Big\}, \qquad \forall i \in \{1, 2, \dots, k\}.
  \label{eq:bucket-formal-def}
\end{equation}
\end{definition}

The resulting partition topology is illustrated in \autoref{fig:k-way-partitioning-model}.

\begin{figure}[H]
  \centering
  \resizebox{0.95\textwidth}{!}{%
  \begin{tikzpicture}[
    scale=1.0,
    bucket/.style={rectangle, draw=blue!80!black, thick, minimum height=1.2cm, align=center, font=\small\bfseries},
    sentinel/.style={font=\footnotesize\bfseries, text=red!80!black},
    pivot/.style={font=\footnotesize\bfseries, text=blue!85!black}
  ]
    % Buckets B_1, B_2, dots, B_k
    \node[bucket, fill=blue!10, minimum width=2.6cm] (b1) at (0, 0) {$B_1$\\[2pt]\footnotesize $s_0 < x \le s_1$};
    \node[bucket, fill=teal!12, minimum width=2.6cm, right=0pt of b1] (b2) {$B_2$\\[2pt]\footnotesize $s_1 < x \le s_2$};
    \node[bucket, fill=gray!12, minimum width=2.4cm, right=0pt of b2] (dots) {$\dots$\\[2pt]\footnotesize $\dots$};
    \node[bucket, fill=orange!12, minimum width=2.6cm, right=0pt of dots] (bk) {$B_k$\\[2pt]\footnotesize $s_{k-1} < x \le s_k$};

    % Delimiters and Pivots
    \node[sentinel, above=6pt of b1.north west] (s0) {$s_0 = -\infty$};
    \node[pivot, above=6pt of b1.north east] (s1) {$s_1$};
    \node[pivot, above=6pt of b2.north east] (s2) {$s_2$};
    \node[pivot, above=6pt of bk.north west] (skm1) {$s_{k-1}$};
    \node[sentinel, above=6pt of bk.north east] (sk) {$s_k = +\infty$};

    % Vertical demarcation lines
    \draw[thick, red!70!black] (b1.north west) -- (b1.south west);
    \draw[thick, blue!80!black] (b1.north east) -- (b1.south east);
    \draw[thick, blue!80!black] (b2.north east) -- (b2.south east);
    \draw[thick, blue!80!black] (bk.north west) -- (bk.south west);
    \draw[thick, red!70!black] (bk.north east) -- (bk.south east);
  \end{tikzpicture}%
  }
  \caption{Partitioning of dataset $S$ into $k$ buckets using $k-1$ sorted pivots.}
  \label{fig:k-way-partitioning-model}
\end{figure}

\subsection{Optimal Balancing Condition}
For the sorting algorithm to maintain optimal I/O complexity, bucket sizes must not deviate excessively from the ideal average size $\frac{n}{k}$.

\begin{property}[Formal Balancing Objective]
All buckets must have asymptotically uniform cardinality:
\begin{equation}
  |B_i| = \Theta\left(\frac{n}{k}\right), \qquad \forall i \in \{1, 2, \dots, k\}.
\end{equation}
\end{property}

The architectural significance of this balance is two-fold:
\begin{enumerate}
  \item \textbf{Asymptotic Optimality of Recursion Depth:} When all buckets remain balanced across recursion levels, tree depth is bounded by:
  \begin{equation}
    \mathcal{O}\left(\log_k \frac{n}{B}\right) = \mathcal{O}\left(\frac{\log(n/B)}{\log(M/B)}\right).
  \end{equation}
  Distributing elements among $k$ buckets at each level requires a single linear scan taking $\mathcal{O}\left(\frac{n}{B}\right)$ I/Os (allocating one output write buffer for each of the $k = \mathcal{O}(M/B)$ buckets). Total I/O complexity matches Aggarwal and Vitter's lower bound:
  \begin{equation}
    \mathcal{O}\left(\frac{n}{B} \log_{M/B} \frac{n}{B}\right) \text{ I/Os}.
  \end{equation}
  \item \textbf{Prevention of Memory Overflow:} No bucket exceeds internal memory capacity during the base case where individual buckets are loaded into RAM and sorted locally.
\end{enumerate}

\FloatBarrier
\section{The Over-Sampling Technique}
\label{sec:oversampling-technique}

\subsection{Shortcomings of Naive Sampling}
Directly extracting only $k - 1$ random pivots from $S$ yields high statistical variance in the spacing between successive pivot keys. Consequently, with high probability, dense clusters and empty gaps emerge, creating pathological buckets with:
\begin{equation}
  |B_i| \gg \frac{n}{k}.
\end{equation}
This severe skew destroys the optimal I/O bounds of external distribution sort.

To sharply suppress this variance and mathematically guarantee that no bucket exceeds a target threshold, we employ \textbf{Over-Sampling}: we draw substantially more than $k$ samples, sort them in RAM, and select the final pivots at equidistant regular strides.

\subsection{Formalization of Over-Sampling Parameters}
We define an over-sampling parameter $a$:
\begin{equation}
  a = \Theta(\log_2 k).
\end{equation}
Specifically, to align analytical constants with exponential concentration bounds, we calibrate:
\begin{equation}
  a + 1 \coloneqq 12 \ln k \qquad (\text{for } k \ge 2).
  \label{eq:oversampling-parameter-a}
\end{equation}

The total size of the random sample pool $A \subset S$ is:
\begin{equation}
  |A| = (a + 1)k - 1.
  \label{eq:sample-size-a}
\end{equation}

\subsection{Pivot Selection Algorithm via Over-Sampling}
The selection workflow proceeds through three sequential phases:
\begin{enumerate}
  \item \textbf{Random Extraction:} Draw $(a + 1)k - 1$ elements uniformly at random from $S$, forming sample set $A$.
  \item \textbf{Sample Sorting:} Sort set $A$ in internal memory (incurring 0 extra disk I/Os since $|A| \le M$):
  \begin{equation}
    A = [A[1], A[2], \dots, A[|A|]], \qquad \text{where } A[1] \le A[2] \le \dots \le A[|A|].
  \end{equation}
  \item \textbf{Equidistant Pivot Selection:} For each index $i \in \{1, 2, \dots, k - 1\}$, pivot $s_i$ is chosen as the element occupying the exact multiple position of $(a + 1)$ in the sorted array $A$:
  \begin{equation}
    s_i \coloneqq A\big[(a + 1) \cdot i\big], \qquad \forall i \in \{1, 2, \dots, k - 1\}.
    \label{eq:pivot-selection-rule}
  \end{equation}
\end{enumerate}

\subsubsection{Structural Layout of the Sorted Sample Vector \texorpdfstring{$A$}{A}}
The sorted array $A$ is partitioned into $k$ logical segments containing $a$ items each, separated by sentinel pivots:
\begin{itemize}
  \item Exactly $a$ samples of $A$ lie strictly between $s_{i-1}$ and $s_i$.
  \item Including the right boundary $s_i$, exactly $a + 1$ samples of $A$ fall into the half-open interval $(s_{i-1}, s_i]$.
\end{itemize}

\begin{figure}[H]
  \centering
  \resizebox{0.95\textwidth}{!}{%
  \begin{tikzpicture}[
    scale=0.95,
    blk/.style={rectangle, draw=blue!80!black, fill=blue!10, thick, minimum height=1.1cm, align=center, font=\small},
    pvt/.style={rectangle, draw=teal!85!black, fill=teal!25, line width=1.4pt, minimum height=1.1cm, minimum width=1.1cm, align=center, font=\small\bfseries}
  ]
    % Blocks and pivot nodes
    \node[blk, minimum width=2.4cm] (b1) at (0, 0) {$a$ items\\ \scriptsize pos. $1 \dots a$};
    \node[pvt, right=0pt of b1] (s1) {$s_1$\\ \tiny $a+1$};
    \node[blk, minimum width=2.6cm, right=0pt of s1] (b2) {$a$ items\\ \scriptsize pos. $a+2 \dots 2a+1$};
    \node[pvt, right=0pt of b2] (s2) {$s_2$\\ \tiny $2(a+1)$};
    \node[blk, fill=gray!12, minimum width=1.8cm, right=0pt of s2] (dots) {$\dots$};
    \node[pvt, right=0pt of dots] (skm1) {$s_{k-1}$\\ \tiny $(k-1)(a+1)$};
    \node[blk, minimum width=2.4cm, right=0pt of skm1] (bk) {$a$ items\\ \scriptsize pos. $(k-1)(a+1)+1 \dots |A|$};

    \node[above=8pt of s1, font=\footnotesize\bfseries, text=teal!90!black] {Pivot 1};
    \node[above=8pt of s2, font=\footnotesize\bfseries, text=teal!90!black] {Pivot 2};
    \node[above=8pt of skm1, font=\footnotesize\bfseries, text=teal!90!black] {Pivot $k-1$};
  \end{tikzpicture}%
  }
  \caption{Structural segmentation of sorted sample array $A$ into $k$ blocks of size $a$ spaced by pivots.}
  \label{fig:oversampling-vector-structure}
\end{figure}

\begin{algorithm}[H]
  \caption{Pivot Selection via Over-Sampling}
  \label{alg:oversampling-pivot-selection}
  \KwInput{Unsorted dataset $S[1 \dots n]$, fan-out $k \ge 2$, parameter $a + 1 = 12 \ln k$}
  \KwOutput{Sorted pivot vector $s = \langle s_1, s_2, \dots, s_{k-1} \rangle$}
  $m \leftarrow (a + 1)k - 1$\;
  $A \leftarrow \operatorname{RandomSample}(S, m)$\tcp*{Draw m items uniformly at random}
  $\operatorname{Sort}(A)$\tcp*{Internal sort in RAM: A[1] <= ... <= A[m]}
  \For{$i \leftarrow 1$ \KwTo $k - 1$}{
    $s_i \leftarrow A[(a + 1) \cdot i]$\tcp*{Equidistant pivots at stride a+1}
  }
  \Return $s$\;
\end{algorithm}

\FloatBarrier
\section{Probabilistic Failure Analysis and Chernoff Bound}
\label{sec:failure-analysis-chernoff}

\subsection{Formal Definition of Failure Event}
Our algorithmic goal is to prevent any bucket from growing excessively large. We establish the critical failure threshold at twice the average bucket capacity $\frac{2n}{k}$.

\begin{definition}[Failure Event ($\text{Failure}$)]
Failure occurs if at least one of the $k$ buckets exceeds the critical threshold:
\begin{equation}
  \text{Failure} \iff \exists i \in \{1, 2, \dots, k\} \text{ such that } |B_i| \ge \frac{2n}{k}.
  \label{eq:failure-definition}
\end{equation}
\end{definition}

We aim to prove that the failure probability is bounded by a constant strictly less than 1:
\begin{equation}
  \mathbb{P}(\text{Failure}) \le \frac{1}{2}.
\end{equation}

\begin{property}[Practical Relevance of $\mathbb{P}(\text{Failure}) \le \frac{1}{2}$]
A failure probability $\le \frac{1}{2}$ is ideal in the analysis of \textbf{Las Vegas} randomized algorithms.
If the algorithm verifies bucket balance a posteriori (costing a linear $\frac{n}{B}$ I/O scan) and restarts sampling upon failure, the expected number of iterations is geometrically distributed with mean:
\begin{equation}
  \mathbb{E}[\text{trials}] = \frac{1}{1 - \mathbb{P}(\text{Failure})} \le \frac{1}{1 - 1/2} = 2.
\end{equation}
Furthermore, after $r$ independent attempts, the probability of remaining unbalanced decays exponentially to $\le \left(\frac{1}{2}\right)^r$, ensuring extreme operational robustness.
\end{property}

\subsection{Ideal Sample Space \texorpdfstring{$S_{\text{sort}}$}{S\_sort} and Fixed-Block Partitioning}
To analyze this stochastic behavior, consider the ideal fully sorted sequence of all $n$ elements in $S$, denoted $S_{\text{sort}}$:
\begin{equation}
  S_{\text{sort}} = [x_1, x_2, \dots, x_n], \qquad \text{with } x_1 \le x_2 \le \dots \le x_n.
\end{equation}
Because every bucket $B_i$ gathers all elements in the interval $(s_{i-1}, s_i]$, the contents of each bucket $B_i$ occupy a contiguous subarray inside $S_{\text{sort}}$.

We conceptually partition the sorted sequence $S_{\text{sort}}$ into fixed, disjoint blocks $t_j$ of uniform size $\frac{2n}{k}$. The total count of such fixed blocks is:
\begin{equation}
  \text{Number of blocks } t_j = \frac{n}{\frac{2n}{k}} = \frac{k}{2}.
\end{equation}
We denote these $\frac{k}{2}$ blocks as $t_1, t_2, \dots, t_{k/2}$.

\begin{figure}[H]
  \centering
  \resizebox{0.95\textwidth}{!}{%
  \begin{tikzpicture}[
    scale=1.0,
    box/.style={rectangle, draw=blue!85!black, fill=blue!10, thick, minimum height=1.1cm, align=center, font=\small},
    brace/.style={decorate, decoration={brace, amplitude=5pt, raise=3pt}, thick}
  ]
    % Blocks t_j of S_sort
    \node[box, minimum width=2.8cm] (t1) at (0, 0) {$t_1$};
    \node[box, minimum width=2.8cm, right=0pt of t1] (t2) {$t_2$};
    \node[box, minimum width=2.8cm, right=0pt of t2] (t3) {$t3$};
    \node[box, fill=gray!12, minimum width=2.0cm, right=0pt of t3] (dots) {$\dots$};
    \node[box, minimum width=2.8cm, right=0pt of dots] (tk2) {$t_{k/2}$};

    % Block width indicators
    \draw[brace, decoration={mirror}] ($(t1.south west)+(0,-0.1)$) -- ($(t1.south east)+(0,-0.1)$)
      node[midway, below=8pt, font=\footnotesize\bfseries, text=blue!85!black] {$\frac{2n}{k}$};
    \draw[brace, decoration={mirror}] ($(t2.south west)+(0,-0.1)$) -- ($(t2.south east)+(0,-0.1)$)
      node[midway, below=8pt, font=\footnotesize\bfseries, text=blue!85!black] {$\frac{2n}{k}$};

    % Full S_sort brace
    \draw[brace, draw=teal!80!black] ($(t1.north west)+(0,0.1)$) -- ($(tk2.north east)+(0,0.1)$)
      node[midway, above=8pt, font=\small\bfseries, text=teal!85!black] {
        $S_{\text{sort}}$ ($n$ items partitioned into $\frac{k}{2}$ contiguous fixed blocks)
      };
  \end{tikzpicture}%
  }
  \caption{Theoretical partitioning of $S_{\text{sort}}$ into $\frac{k}{2}$ fixed blocks $t_j$ of size $\frac{2n}{k}$.}
  \label{fig:fixed-blocks-partition}
\end{figure}

\subsection{Logical Implication of the Failure Event}
Examine the geometric configuration when failure occurs for an arbitrary bucket $B_i$, meaning $|B_i| \ge \frac{2n}{k}$:
\begin{enumerate}
  \item Since $B_i$ forms a contiguous interval in $S_{\text{sort}}$ of length at least $\frac{2n}{k}$, it must completely contain (or span across) at least one fixed block $t_j$:
  \begin{equation}
    \exists j \in \left\{1, \dots, \frac{k}{2}\right\} \text{ such that } t_j \subseteq B_i.
  \end{equation}
  \item The delimiting pivots $s_{i-1}$ and $s_i$ sit strictly outside block $t_j$ (to the left and right, respectively).
  \item By construction of over-sampling, the entire bucket $B_i$ contains at most $a$ interior samples, plus boundary pivot $s_i$. Thus, the total count of samples from $A$ falling within $B_i$ is at most $a + 1$.
  \item Since $t_j \subseteq B_i$ and pivot $s_i$ lies outside $t_j$, fixed block $t_j$ can contain strictly fewer than $a + 1$ samples of $A$:
  \begin{equation}
    |t_j \cap A| < a + 1.
  \end{equation}
\end{enumerate}

Thus, the logical implication holds:
\begin{equation}
  X \implies Y,
\end{equation}
where:
\begin{align}
  X &\colon \exists B_i \text{ such that } |B_i| \ge \frac{2n}{k}, \\
  Y &\colon \exists t_j \text{ such that } t_j \text{ contains } < a + 1 \text{ samples of } A.
\end{align}

By fundamental probability laws:
\begin{equation}
  X \implies Y \implies \mathbb{P}(X) \le \mathbb{P}(Y).
\end{equation}
We upper-bound the failure probability by analyzing the likelihood that at least one fixed block receives too few samples:
\begin{equation}
  \mathbb{P}(\text{Failure}) \le \mathbb{P}\left(\exists j \in \left\{1, \dots, \frac{k}{2}\right\} \colon |t_j \cap A| < a + 1\right).
\end{equation}

\subsection{Union Bound}
Applying Boole's inequality (Union Bound) across all $\frac{k}{2}$ fixed blocks:
\begin{equation}
  \mathbb{P}\left(\bigcup_{j=1}^{k/2} \big\{|t_j \cap A| < a + 1\big\}\right) \le \sum_{j=1}^{k/2} \mathbb{P}\big(|t_j \cap A| < a + 1\big).
\end{equation}

Because each block $t_j$ has identical length $|t_j| = \frac{2n}{k}$, the probability of an individual sample landing in $t_j$ is identical for every $j$. Selecting an arbitrary reference block $t^*$:
\begin{equation}
  \mathbb{P}(\text{Failure}) \le \frac{k}{2} \cdot \mathbb{P}\big(t^* \text{ contains } < a + 1 \text{ samples of } A\big).
  \label{eq:union-bound-failure}
\end{equation}

\subsection{Expectation Computation}
Fix block $t^*$. Let $X$ be the random variable counting how many sample items from $A$ land in $t^*$. Define indicator variables for each drawn sample $r \in \{1, \dots, |A|\}$:
\begin{equation}
  X_r \coloneqq \begin{cases}
    1 & \text{if sample } r \text{ of } A \text{ belongs to } t^*, \\
    0 & \text{otherwise.}
  \end{cases}
\end{equation}
Total random count is:
\begin{equation}
  X = \sum_{r=1}^{|A|} X_r.
\end{equation}

The marginal probability that an arbitrary uniform sample from $S$ falls into $t^*$ is:
\begin{equation}
  \mathbb{P}(\text{sample } \in t^*) = \frac{|t^*|}{n} = \frac{\frac{2n}{k}}{n} = \frac{2}{k}.
\end{equation}

By linearity of expectation:
\begin{equation}
  \mathbb{E}[X] = \mathbb{E}\left[ \sum_{r=1}^{|A|} X_r \right] = \sum_{r=1}^{|A|} \mathbb{E}[X_r] = \sum_{r=1}^{|A|} \mathbb{P}(X_r = 1) = |A| \cdot \frac{2}{k}.
\end{equation}

Substituting total sample size $|A| = (a + 1)k - 1$:
\begin{align}
  \mathbb{E}[X] &= \big[(a + 1)k - 1\big] \cdot \frac{2}{k} \notag\\
  &= 2(a + 1) - \frac{2}{k}.
\end{align}

Since $k \ge 2$, the subtracted fraction satisfies:
\begin{equation}
  \frac{2}{k} \le \frac{2}{2} = 1.
\end{equation}
We lower-bound the expectation:
\begin{equation}
  \mathbb{E}[X] \ge 2(a + 1) - 1.
\end{equation}

Since $a \ge 1$ (as $a + 1 = 12 \ln k \ge 12 \ln 2 \approx 8.31$ for all $k \ge 2$):
\begin{align}
  2(a + 1) - 1 &= \frac{3}{2}(a + 1) + \underbrace{\frac{1}{2}(a + 1) - 1}_{\ge 0} \notag\\
  &\ge \frac{3}{2}(a + 1).
\end{align}
We establish the key inequality relating expectation to threshold:
\begin{equation}
  \mathbb{E}[X] \ge \frac{3}{2}(a + 1) \iff a + 1 \le \frac{2}{3}\mathbb{E}[X].
  \label{eq:expectation-lower-bound}
\end{equation}

\subsection{Chernoff Bound Application (Lower Tail)}
The Chernoff lower-tail concentration bound establishes that for a sum $X = \sum X_r$ of independent random variables bounded in $[0, 1]$, for any relative deviation $\delta \in (0, 1)$:
\begin{equation}
  \mathbb{P}\big(X < (1 - \delta)\mathbb{E}[X]\big) \le e^{-\frac{\delta^2}{2}\mathbb{E}[X]}.
  \label{eq:chernoff-lower-tail-theorem}
\end{equation}

\begin{property}[Intuitive Interpretation of Chernoff Bound]
Random variable $X$ concentrates sharply around its mean $\mathbb{E}[X]$. The probability of falling far below the expectation decays exponentially with respect to deviation square $\delta^2$.
\end{property}

We match $\mathbb{P}(X < a + 1)$ with Chernoff's format. From inequality~\eqref{eq:expectation-lower-bound}, $a + 1 \le \frac{2}{3}\mathbb{E}[X]$. Setting relative deviation:
\begin{equation}
  1 - \delta = \frac{2}{3} \implies \delta = 1 - \frac{2}{3} = \frac{1}{3}.
\end{equation}

Consequently:
\begin{equation}
  \mathbb{P}(X < a + 1) \le \mathbb{P}\left(X < \frac{2}{3}\mathbb{E}[X]\right) = \mathbb{P}\left(X < \left(1 - \frac{1}{3}\right)\mathbb{E}[X]\right).
\end{equation}

Applying Chernoff bound~\eqref{eq:chernoff-lower-tail-theorem} with $\delta = \frac{1}{3}$:
\begin{equation}
  \mathbb{P}\left(X < \left(1 - \frac{1}{3}\right)\mathbb{E}[X]\right) \le e^{-\frac{(1/3)^2}{2}\mathbb{E}[X]} = e^{-\frac{1/9}{2}\mathbb{E}[X]} = e^{-\frac{1}{18}\mathbb{E}[X]}.
\end{equation}

Substituting our expectation lower bound $\mathbb{E}[X] \ge \frac{3}{2}(a + 1)$:
\begin{align}
  e^{-\frac{1}{18}\mathbb{E}[X]} &\le e^{-\frac{1}{18} \cdot \frac{3}{2}(a + 1)} \notag\\
  &= e^{-\frac{3}{36}(a + 1)} \notag\\
  &= e^{-\frac{1}{12}(a + 1)}.
\end{align}

\subsection{Exact Calibration of \texorpdfstring{$a$}{a} and Conclusion}
Substituting our chosen parameter $a + 1 = 12 \ln k$:
\begin{equation}
  e^{-\frac{1}{12}(a + 1)} = e^{-\frac{1}{12} \cdot (12 \ln k)} = e^{-\ln k} = \frac{1}{k}.
\end{equation}

Therefore, for any individual reference block $t^*$:
\begin{equation}
  \mathbb{P}(t^* \text{ contains } < a + 1 \text{ samples}) \le \frac{1}{k}.
\end{equation}

Multiplying by factor $\frac{k}{2}$ from the Union Bound~\eqref{eq:union-bound-failure}:
\begin{equation}
  \mathbb{P}(\text{Failure}) \le \frac{k}{2} \cdot \mathbb{P}\big(t^* \text{ contains } < a + 1 \text{ samples}\big) \le \frac{k}{2} \cdot \frac{1}{k} = \frac{1}{2}.
\end{equation}

\begin{theorem}[Over-Sampling Balance Theorem]
With success probability at least $1 - \frac{1}{2} = \frac{1}{2}$, over-sampling yields a partition where all buckets $B_i$ strictly satisfy size bound $|B_i| < \frac{2n}{k}$:
\begin{equation}
  \mathbb{P}\left(\forall i \in \{1, \dots, k\} \colon |B_i| < \frac{2n}{k}\right) \ge \frac{1}{2}.
\end{equation}
This establishes uniform balancing and closes the I/O efficiency analysis of External Distribution Sort. \qedhere
\end{theorem}

\FloatBarrier
\section{Introduction to Random Sampling and Non-Streaming Limits}
\label{sec:intro-random-sampling-non-streaming}

Following the analysis of distribution sort, we must study the extraction primitive itself: how to formally draw $m$ random elements from a universe $S$ of size $n$?

\subsection{Formal Problem Definition}
Given an indexed array $S[1 \dots n]$ containing $n$ elements, extract subset $D \subseteq S$ of target cardinality $m$ satisfying:
\begin{itemize}
  \item \textbf{Marginal Uniformity:} Each item $i \in S$ shares identical inclusion probability:
  \begin{equation}
    \mathbb{P}(i \in D) = \frac{m}{n}.
  \end{equation}
  \item \textbf{Exact Cardinality Without Replacement:} The sample contains exactly $m$ distinct items:
  \begin{equation}
    |D| = m.
  \end{equation}
\end{itemize}

We assume access to an integer random generator primitive:
\begin{equation}
  \operatorname{Rand}(a, b) \to r \in [a, b], \qquad \text{where } r \text{ is drawn uniformly.}
\end{equation}

\subsection{Algorithm 1: Rejection Sampling (Naive Insertion Sampling)}
The simplest scheme draws random candidate indices, inserting new ones into a dynamic set while discarding duplicates until reaching size $m$.

\begin{algorithm}[H]
  \caption{Rejection Sampling via Dynamic Set Insertion}
  \label{alg:rejection-sampling}
  \KwInput{Index universe $[1 \dots n]$, target sample size $m \le n$}
  \KwOutput{Sorted set $D$ containing $m$ distinct indices}
  $D \leftarrow \emptyset$\;
  \While{$|D| < m$}{
    $p \leftarrow \operatorname{Rand}(1, n)$\;
    \If{$p \notin D$}{
      $D \leftarrow D \cup \{p\}$\;
    }
  }
  $\operatorname{Sort}(D)$\;
  \Return $D$\;
\end{algorithm}

\subsubsection{Collision Analysis and Failure Probability}
At each step of the \texttt{while} loop, the probability that candidate $p$ is already in $D$ (causing a collision and a wasted attempt) is:
\begin{equation}
  \mathbb{P}(\text{collision}) = \mathbb{P}(p \in D) = \frac{|D|}{n}.
\end{equation}

Across execution, $|D|$ rises from $0$ to $m - 1$. The collision probability at any step is bounded by:
\begin{equation}
  \mathbb{P}(\text{failure}) < \frac{m}{n}.
\end{equation}

\begin{property}[Operating Regimes]
Efficiency depends heavily on the ratio of $m$ to $n$:
\begin{itemize}
  \item \textbf{Sparse Regime ($m \le \frac{n}{2}$):} Step failure probability satisfies $\mathbb{P}(\text{failure}) < \frac{1}{2}$. Expected draws before success is less than 2, ensuring fast execution. Smaller $m/n$ ratios yield even fewer collisions.
  \item \textbf{Dense Regime ($m \to n$):} As $m$ nears $n$, drawing an unselected element becomes increasingly improbable, degenerating into the classical \textbf{Coupon Collector's Problem} and inflating expected iterations to:
  \begin{equation}
    \Theta(n \ln m).
  \end{equation}
\end{itemize}
\end{property}

\subsubsection{Impact of Underlying Data Structures}
As highlighted in lecture:
\begin{itemize}
  \item \textbf{Set Management for $D$:} Implementing $D$ as a \textbf{hash table} enables membership checks $p \notin D$ and insertions in expected $\mathcal{O}(1)$ time. An unsorted array would require linear search taking $\mathcal{O}(|D|) = \mathcal{O}(m)$ per draw.
  \item \textbf{Final Sorting:} Operation $\operatorname{Sort}(D)$ consumes $\mathcal{O}(m \log m)$ CPU time to return sorted positions.
\end{itemize}

\subsection{Handling Variable-Length Records (Strings)}
How do we sample when elements of $S$ are arbitrary variable-length records or text strings?

\begin{figure}[H]
  \centering
  \resizebox{0.95\textwidth}{!}{%
  \begin{tikzpicture}[
    scale=0.95,
    str/.style={rectangle, draw=blue!85!black, fill=blue!10, thick, minimum height=1.0cm, align=center, font=\small\ttfamily},
    ptr/.style={rectangle, draw=teal!85!black, fill=teal!15, thick, minimum height=0.9cm, minimum width=1.1cm, align=center, font=\footnotesize\bfseries},
    arrow/.style={-{Stealth[scale=0.85]}, thick, draw=teal!90!black}
  ]
    % Pointer array
    \node[ptr] (p1) at (0, 2) {$P[1]$};
    \node[ptr, right=0pt of p1] (p2) {$P[2]$};
    \node[ptr, right=0pt of p2] (p3) {$P[3]$};
    \node[ptr, right=0pt of p3] (p4) {$P[4]$};
    \node[right=8pt of p4, font=\small\bfseries, text=teal!90!black] {Indexed contiguous pointer array ($1 \dots n$)};

    % Variable length string records
    \node[str, minimum width=1.8cm] (s1) at (0, 0) {"data"};
    \node[str, minimum width=3.2cm, right=0pt of s1] (s2) {"algorithms"};
    \node[str, minimum width=2.4cm, right=0pt of s2] (s3) {"system"};
    \node[str, minimum width=1.4cm, right=0pt of s3] (s4) {"io"};
    \node[right=8pt of s4, font=\small\bfseries, text=blue!85!black] {Variable-length disk records};

    % Pointer arrows
    \draw[arrow] (p1.south) -- (s1.north);
    \draw[arrow] (p2.south) -- (s2.north);
    \draw[arrow] (p3.south) -- (s3.north);
    \draw[arrow] (p4.south) -- (s4.north);
  \end{tikzpicture}%
  }
  \caption{Mapping variable-length string records using a contiguous indexed pointer array.}
  \label{fig:variable-length-pointers}
\end{figure}

In a raw sequential file of variable-length strings, scalar index $p \in [1, n]$ from $\operatorname{Rand}(1, n)$ has no immediate $\mathcal{O}(1)$ byte offset.

To solve this, we construct a contiguous \textbf{array of pointers} $P[1 \dots n]$, where pointer $P[i]$ references the exact byte offset of the $i$-th record. Random draws operate directly over pointer indices in $\mathcal{O}(1)$ time, granting instant access to the targeted record.

\subsection{Why Algorithm 1 Fails in the Streaming Model}
The lecture concludes with a vital structural observation bridging into data streams:

\begin{property}[Professor's Critical Observation: \textit{``Our algorithm is not Streaming!''}]
Algorithm~\ref{alg:rejection-sampling} suffers from fundamental constraints preventing its application to data streams:
\begin{enumerate}
  \item \textbf{Arbitrary Random Access:} It requires querying any arbitrary index $p \in [1, n]$ on demand at any iteration.
  \item \textbf{Violation of Single-Pass Constraint:} Stream elements arrive in strict sequential order; rewinding or random hopping is strictly prohibited.
  \item \textbf{Requirement of Known $n$:} Primitive $\operatorname{Rand}(1, n)$ requires knowing total cardinality $n$ upfront before processing begins.
\end{enumerate}
\end{property}

This limitation motivates the development of \textbf{Random Sampling Algorithms on Data Streams} (examined in \autoref{chap:random-sampling-data-streams}), where cardinality $n$ may be unknown and admission decisions must be finalized \emph{on-the-fly} as elements arrive.
